\subsection{Operadores diferenciales sobre un $S$-preesquema diferencialmente
liso}
\label{IV.16.11}

\begin{env}[16.11.1]
\label{IV.16.11.1}
Sean $f:X\to S$ un morfismo, $U$ un abierto de $X$,
$(z_\lambda)_{\lambda\in L}$ una familia de secciones de $\mathcal{O}_X$
sobre $U$ tal que los $dz_\lambda$ formen un sistema de generadores
de $\Omega_{X/S}^1|U=\Omega_{U/S}^1$. Sea $m$ un entero o el símbolo
$\infty$, y pongamos, para todo $\lambda$
\[
  \zeta_\lambda=\delta z_\lambda=d^m z_\lambda-z_\lambda
  \in\Gamma(U,\mathscr{P}_{X/S}^m).
  \tag{16.11.1.1}
  \label{IV.16.11.1.1}
\]

Usaremos también las notaciones usuales del análisis; para todo
$\mathbf{p}=(p_\lambda)\in\mathbb{N}^{(L)}$ (con $p_\lambda=0$ salvo para un
número finito de índices), pondremos
\[
  |\mathbf{p}|=\sum_\lambda p_\lambda,
  \qquad \mathbf{p}!=\prod_\lambda(p_\lambda!)
  \tag{16.11.1.2}
  \label{IV.16.11.1.2}
\]
\[
  \binom{\mathbf{p}}{\mathbf{q}}
  =\mathbf{p}!/(\mathbf{q}!(\mathbf{p}-\mathbf{q})!)
  \quad\text{para $\mathbf{p},\mathbf{q}$ en $\mathbb{N}^{(L)}$,
  $\mathbf{q}\leq\mathbf{p}$}
  \tag{16.11.1.3}
  \label{IV.16.11.1.3}
\]
y convenimos que $\binom{\mathbf{p}}{\mathbf{q}}=0$ si
$\mathbf{q}\nleq\mathbf{p}$,
\[
  \mathbf{z}^{\mathbf{p}}=\prod_\lambda(z_\lambda)^{p_\lambda},
  \qquad
  \boldsymbol{\zeta}^{\mathbf{p}}=\prod_\lambda(\zeta_\lambda)^{p_\lambda}.
  \tag{16.11.1.4}
  \label{IV.16.11.1.4}
\]

Por tanto, con estas notaciones
\[
  d^m(\mathbf{z}^{\mathbf{p}})
  =(d^m(\mathbf{z}))^{\mathbf{p}}
  =(\boldsymbol{\zeta}+\mathbf{z})^{\mathbf{p}}
  =\sum_{\mathbf{q}\leq\mathbf{p}}
    \binom{\mathbf{p}}{\mathbf{q}}
    \mathbf{z}^{\mathbf{p}-\mathbf{q}}\boldsymbol{\zeta}^{\mathbf{q}}
  \tag{16.11.1.5}
  \label{IV.16.11.1.5}
\]
\[
  \boldsymbol{\zeta}^{\mathbf{p}}
  =(d^m\mathbf{z}-\mathbf{z})^{\mathbf{p}}
  =\sum_{\mathbf{q}\leq\mathbf{p}}
    (-1)^{|\mathbf{p}-\mathbf{q}|}
    \binom{\mathbf{p}}{\mathbf{q}}
    \mathbf{z}^{\mathbf{p}-\mathbf{q}}d^m(\mathbf{z}^{\mathbf{q}}).
  \tag{16.11.1.6}
  \label{IV.16.11.1.6}
\]

Como los $dz_\lambda$ generan $\Omega_{X/S}^1$ y son las imágenes de
$\delta z_\lambda$, y puesto que el homomorfismo canónico (16.3.1.1) es
sobreyectivo, se concluye que para $m$ finito los $\delta z_\lambda$ generan
la $\mathcal{O}_U$-álgebra $\mathscr{P}_{U/S}^m$ (Bourbaki,
\emph{Alg. comm.}, cap.~III, §~2, n\textsuperscript{o}~8, cor.~2 del teorema~1).
Por tanto, los $\boldsymbol{\zeta}^{\mathbf{p}}$ (para $|\mathbf{p}|\leq m$)
generan el $\mathcal{O}_U$-módulo $\mathscr{P}_{U/S}^m$. Un operador
diferencial $D\in\operatorname{Diff}_{U/S}^m$ queda por consiguiente
completamente determinado por los valores de
$\langle\boldsymbol{\zeta}^{\mathbf{p}},D\rangle$ para
$|\mathbf{p}|\leq m$, o, lo que es equivalente en virtud de
(16.11.1.5) y (16.11.1.6), por los valores
\oldpage[IV-4]{54}
de $\langle d^m(\mathbf{z}^{\mathbf{p}}),D\rangle=D(\mathbf{z}^{\mathbf{p}})$
para $|\mathbf{p}|\leq m$; más precisamente, se deduce de (16.11.1.5) que
se tiene
\[
  D(\mathbf{z}^{\mathbf{p}})
  =\langle d^m(\mathbf{z}^{\mathbf{p}}),D\rangle
  =\sum_{\mathbf{q}\leq\mathbf{p}}
    \binom{\mathbf{p}}{\mathbf{q}}
    \langle\boldsymbol{\zeta}^{\mathbf{q}},D\rangle
    \mathbf{z}^{\mathbf{p}-\mathbf{q}}.
  \tag{16.11.1.7}
  \label{IV.16.11.1.7}
\]
\end{env}

\begin{theorem}[16.11.2]
\label{IV.16.11.2}
Sean $f:X\to S$ un morfismo, $U$ un abierto de $X$,
$(z_\lambda)_{\lambda\in L}$ una familia de secciones de $\mathcal{O}_X$
sobre $U$ tal que la familia $(dz_\lambda)_{\lambda\in L}$ genere
$\Omega_{X/S}^1|U=\Omega_{U/S}^1$. Las condiciones siguientes son
equivalentes:
\begin{enumerate}
  \item[a)] $f|U$ es diferencialmente liso y $(dz_\lambda)$ es una base
  del $\mathcal{O}_U$-módulo $\Omega_{U/S}^1$.
  \item[b)] Existe una familia
  $(D_{\mathbf{p}})_{\mathbf{p}\in\mathbb{N}^{(L)}}$ de operadores
  diferenciales de $\mathcal{O}_U$ en sí mismo que satisface las condiciones
  \[
    D_{\mathbf{p}}(\mathbf{z}^{\mathbf{q}})
    =\binom{\mathbf{q}}{\mathbf{p}}
      \mathbf{z}^{\mathbf{q}-\mathbf{p}}
    \quad(\mathbf{p},\mathbf{q}\text{ en }\mathbb{N}^{(L)}).
    \tag{16.11.2.1}
    \label{IV.16.11.2.1}
  \]
\end{enumerate}
Además, cuando estas condiciones se satisfacen, la familia
$(D_{\mathbf{p}})$ queda determinada de forma única por las condiciones
(16.11.2.1) y satisface las relaciones
\[
  D_{\mathbf{p}}\circ D_{\mathbf{q}}
  =D_{\mathbf{q}}\circ D_{\mathbf{p}}
  =\frac{(\mathbf{p}+\mathbf{q})!}{\mathbf{p}!\mathbf{q}!}
    D_{\mathbf{p}+\mathbf{q}}
  \quad(\mathbf{p},\mathbf{q}\text{ en }\mathbb{N}^{(L)}).
  \tag{16.11.2.2}
  \label{IV.16.11.2.2}
\]
Finalmente, si $L$ es finito, para todo entero $m$, los $D_{\mathbf{p}}$ tales que
$|\mathbf{p}|\leq m$ forman una base del $\mathcal{O}_U$-módulo
$\operatorname{Diff}_{U/S}^m$, es decir, todo operador diferencial
de orden $\leq m$ sobre $U$ se escribe de una sola manera bajo la forma
\[
  D=\sum_{|\mathbf{p}|\leq m}a_{\mathbf{p}}D_{\mathbf{p}}
\]
donde los $a_{\mathbf{p}}$ son secciones de $\mathcal{O}_X$ sobre
$U$.
\end{theorem}

Observemos primero que, en virtud de (16.11.1.6) y (16.11.1.5), se verifica
inmediatamente que las condiciones (16.11.2.1) son equivalentes a
\[
  \langle\boldsymbol{\zeta}^{\mathbf{p}},D_{\mathbf{q}}\rangle
  =\delta_{\mathbf{p}\mathbf{q}}
  \quad\text{(índice de Kronecker)}.
  \tag{16.11.2.3}
  \label{IV.16.11.2.3}
\]
La existencia de la familia $(D_{\mathbf{p}})$ que satisface estas relaciones
implica primero (tomando $|\mathbf{p}|=1$) que los $dz_\lambda$ son
linealmente independientes y, por tanto, forman una \emph{base} del
$\mathcal{O}_U$-módulo $\Omega_{U/S}^1$. Luego, para todo entero $m\geq1$, se
deduce del mismo modo de (16.11.2.3) que los $\boldsymbol{\zeta}^{\mathbf{p}}$
tales que $|\mathbf{p}|\leq m$ son linealmente independientes; por consiguiente
el homomorfismo canónico (16.3.1.1) es inyectivo, luego biyectivo, y esto
demuestra que b) implica a). La recíproca se deduce inmediatamente de la
definición (16.10.1), pues el hecho de que los $\boldsymbol{\zeta}^{\mathbf{p}}$
formen una base de $\mathscr{P}_{U/S}^m$ para $|\mathbf{p}|\leq m$ implica la
existencia y unicidad de una familia de homomorfismos
$u_{\mathbf{q},m}:\mathscr{P}_{U/S}^m\to\mathcal{O}_U$
($|\mathbf{q}|\leq m$) tales que
$\langle\boldsymbol{\zeta}^{\mathbf{p}},u_{\mathbf{q},m}\rangle
=\delta_{\mathbf{p}\mathbf{q}}$ para $|\mathbf{p}|\leq m$,
$|\mathbf{q}|\leq m$. Para un valor dado de $\mathbf{q}$, los operadores
diferenciales correspondientes a los $u_{\mathbf{q},m}$ para
$m\geq|\mathbf{q}|$ se identifican con un mismo operador $D_{\mathbf{q}}$.
Esto demuestra, pues, que a) implica b), y además que la familia
$(D_{\mathbf{p}})$ queda determinada de forma única y que, si $L$ es finito,
para $|\mathbf{p}|\leq m$, los $D_{\mathbf{p}}$ forman una base del dual
$\operatorname{Diff}_{U/S}^m$ de $\mathscr{P}_{U/S}^m$. Finalmente, las relaciones
(16.11.2.2) se deducen inmediatamente de la expresión de los valores de los tres
operadores considerados sobre los $\mathbf{z}^{\mathbf{r}}$, y del hecho de que
los $\boldsymbol{\zeta}^{\mathbf{r}}$ para $|\mathbf{r}|\leq m$ generan
$\mathscr{P}_{U/S}^m$.

\oldpage[IV-4]{55}

\begin{remarks}[16.11.3]
\label{IV.16.11.3}
\begin{enumerate}
  \item[(i)] El hecho de que los $D_{\mathbf{p}}$ sean permutables dos a dos en
  virtud de (16.11.2.2) no implica naturalmente que la
  $\mathcal{O}_U$-álgebra $\operatorname{Diff}_{U/S}$ sea conmutativa, pues los
  $D_{\mathbf{p}}$ no conmutan con la multiplicación por secciones de
  $\mathcal{O}_U$ salvo si $\mathbf{p}=0$.\footnote{El texto impreso dice
  $n=0$; debe ser $\mathbf{p}=0$, pues $D_{\mathbf{p}}$ es de orden
  $|\mathbf{p}|$.}
  \item[(ii)] Los índices $\mathbf{p}$ tales que $|\mathbf{p}|=1$ son los
  $\boldsymbol{\varepsilon}_\lambda
  =(\varepsilon_{\lambda\mu})_{\mu\in L}$ con
  $\varepsilon_{\lambda\mu}=0$ si $\mu\neq\lambda$ y
  $\varepsilon_{\lambda\lambda}=1$; cuando $L$ es finito, los operadores
  $D_{\boldsymbol{\varepsilon}_\lambda}$ no son otros que las
  $S$-derivaciones $D_i$ introducidas en (16.5.7). Obsérvese que, en general
  (y contrariamente a lo que ocurre en el análisis clásico), no es cierto que
  un operador diferencial de orden cualquiera pueda escribirse como combinación
  lineal de potencias de los $D_i$ (cf. (16.12)).
  \item[(iii)] Para todo entero $r\geq1$, se puede definir la noción de
  morfismo \emph{diferencialmente liso hasta el orden $r$} sustituyendo
  en (16.10.1) la condición (ii) por la condición de que los homomorfismos
  \[
    \mathbf{S}_{\mathcal{O}_X}^m(\Omega_{X/S}^1)
    \longrightarrow\mathscr{G}r_m(\mathscr{P}_{X/S})
  \]
  sean biyectivos \emph{para todo $m\leq r$}. El razonamiento de (16.11.2)
  demuestra entonces que si, en la condición a), se sustituye «diferencialmente
  liso» por «diferencialmente liso hasta el orden $r$», esta condición es
  equivalente a la condición b) en la que nos restringimos a los
  $\mathbf{p}\in\mathbb{N}^{(L)}$, $\mathbf{q}\in\mathbb{N}^{(L)}$ tales que
  $|\mathbf{p}|\leq r$, $|\mathbf{q}|\leq r$.
\end{enumerate}
\end{remarks}
